\documentclass{article}
\usepackage[margin=2cm]{geometry}
\usepackage{amsmath}
\usepackage{graphicx}
\usepackage{booktabs}
\usepackage[utf8]{inputenc}
\begin{document}
\subsection*{0.1 Infinite countertarms}
We now turn towards the identification of the infinite countertarms in the more general curved domain wall case. We may first solve for all of the infinite countertarms via the usual holographic renormalization procedure. Once we have these, we will
1. Check that in the flat limit, they reduce to the divergent pieces of the flat countertarms found above.
2. Add to them the finite pieces found in but missing in the holographic renormalization procedure.
For simplicity, we will perform holographic renormalization on supersymmetric solutions only, and thus the infinite countertarms we obtain are universal for supersymmetric solutions only. We begin by using the expression for the on-shell Ricci scalar,
\[
R=4\left(\sigma^{\prime}\right)^{2}+\left[-\left(\phi^{3 \prime}\right)^{2}+\cos^{2}\phi^{3}\left(\phi^{0 \prime}\right)^{2}\right]+6 V
\]
to rewrite the action as
\[
S_{6 \mathrm{D}}=-\frac{1}{2} \int d u d^{5} x \sqrt{g} \, e^{5 f} V
\]
We have not included the Gibbons-Hawking term yet, but will do so later. The first step of holographic renormalization is to isolate the divergent terms. We may do so by expanding all fields using their UV asymptotics, then integrating over small \(z\) and evaluating on the cutoff \(\epsilon\). Doing so, we find
\[
\begin{aligned}
S_{6 \mathrm{D}}= & -\frac{1}{2} \int d^{5} x \sqrt{g} e^{5 f_{k}}\left[\frac{1}{\epsilon^{5}}+\frac{1}{3 \epsilon^{3}}\left(25 f_{2}+\left(\phi_{1}^{0}\right)^{2}\right)\right. \\
& \left.+\frac{1}{24 \epsilon}\left(1500 f_{2}^{2}+600 f_{4}+120 f_{2}\left(\phi_{1}^{0}\right)^{2}-\left(\phi_{1}^{0}\right)^{4}\right. \\
& \left.+48 \phi_{1}^{0} \phi_{3}^{0}+36\left(-\left(\phi_{2}^{3}\right)^{2}+4 \sigma_{2}^{2}\right)\right)
\end{aligned}
\]
where we've thrown out all non-divergent contributions. Note that the integration would naively give a \(\log \epsilon\), but this vanishes on the BPS equations since they constrain the UV asymptotic expansion coefficients in the following way,
\[
25 f_{5}+2 \phi_{1}^{0} \phi_{4}^{0}-3 \phi_{2}^{3} \phi_{3}^{3}+12 \sigma_{2} \sigma_{3}=0
\]
The absence of the logarithmic term is to be expected, since any dual five-dimensional field theory is anomaly-free. The Gibbons-Hawking term is
\[
S_{\mathrm{GH}}=-\frac{5}{2} \int d^{5} x \sqrt{g} \, e^{5 f} f^{\prime}
\]
\end{document}
